\section{Construction of the Liquidation Measure}
\label{app:data}

The liquidation series is extracted from Dune Analytics query
\texttt{6912877}, versioned in the replication package, over the
open-source Spellbook tables \texttt{lending.borrow} (debt leg) and
\texttt{lending.supply} (collateral leg), across all EVM chains and lending
protocols indexed at extraction time. This appendix records the
construction decisions that materially affect the primary regressor. The
full query is in the package.

\paragraph{Pre-aggregation before the join.}
An event-level join of the debt and collateral legs on transaction produces
a Cartesian product for transactions containing multiple debt and multiple
collateral events. This pattern is common in bundled liquidation bots, and it
inflates totals. The query pre-aggregates each leg by transaction before
joining, eliminating the inflation by construction. The correction revised
the 2024 and 2025 aggregate volumes downward by roughly fourteen and thirty
percent respectively relative to a naive join, while pre-2024 aggregates
moved by less than a quarter of a percent.

\paragraph{Primary measure.}
The primary variable is the hourly sum of debt forcibly repaid in USD,
zero-filled on hours without qualifying liquidations, entered as
$L_{t-1} = \ln(1 + \text{liq}^{\text{USD}}_{t-1})$ with a one-hour lag. The
seized-collateral leg is retained as a robustness alternative. The two
measures correlate above 0.97, and the choice changes the headline
coefficient by less than three percent in log space.

\paragraph{Collateral filter.}
Only liquidations seizing ETH-like collateral are kept, using the symbol set
\{WETH, ETH, stETH, wstETH, rETH, cbETH, sfrxETH, ETHx\}, so the measure
isolates the ETH leverage channel. Liquid-staking derivatives are included,
so the collateral is ETH-like in the broad sense. A depeg of a staking
derivative could in principle inject non-pure-ETH liquidations, a rare
event we note for completeness.

\paragraph{Exclusions.}
Three exclusions are applied on top of the structural filters. Positions
with absolute debt below five dollars are dropped, removing a cluster of
roughly 75{,}000 micro-events in September 2025 linked to bot activity that
is immaterial in dollar terms. UwuLend liquidations from June 2024 onward
are dropped to exclude an exploit-driven hour reporting an implausible
volume attributable to on-chain exploit movements combined with pricing
errors on illiquid tokens. Euler liquidations during the March--April 2023
hack window are dropped; Euler v2, a distinct protocol launched in 2024, is
retained. A small number of hours in late 2025 carry anomalous
collateral-to-debt ratios from oracle pricing of illiquid tokens; these are
economically invisible in the debt-repaid measure.

\paragraph{Coverage caveats.}
MakerDAO, historically the largest on-chain liquidator, sits outside the
Spellbook lending schema (a CDP architecture) and is de facto absent, so the
measure understates total on-chain deleveraging, most severely in 2021. No
centralized-exchange liquidations enter. The query was executed once, on a
single date; a future re-execution may differ marginally if the Spellbook
tables are updated.

\paragraph{Sensitivity to the MakerDAO gap.}
Because the MakerDAO omission is the largest coverage gap, we bound its
effect by an actual re-fit rather than an assumption. We inflate the observed
liquidation series by a missing-Maker share $m$ (up to a punitive $m=0.5$)
and re-estimate the locked specification, placing the missing mass three
ways: proportional to currently-active hours, on stress-adjacent zero hours,
and weighted toward high open-interest (leverage) hours. The shock
coefficient's sign is preserved under every allocation and horizon. At impact
its magnitude moves by at most \MakerLevelImpactPct\%, but at twelve hours the
leverage-weighted infill (the largest-displacement allocation, since Maker
auctions cluster in high-leverage, falling-price hours) attenuates the
deepening by up to \MakerLevelMediumPct\%. Both displacement figures are the
maximum over the defensible allocation range $m \le 0.25$; the punitive
$m=0.5$ end of the sweep changes neither sign nor significance. The
medium-horizon magnitude is therefore coverage-sensitive. The bounded object is not: the per-period exceedance
asymmetry at the 1\% level stays within \MakerExcOverSesoi{} of the smallest
effect size of interest under every allocation, and shrinks toward zero under
the leverage-weighted and stress-adjacent infills, so a more complete measure
would tighten rather than overturn the null.